% element: 10-s011:e05
\BeleskeID{10-s011}
Karakteristiku $V_0=f(V_I)$ određujemo promenom ulaznog napona od nule do $V_{CC}$, uz neopterećen izlaz $I_O=0$. Krećemo od $V_I=0$ i proveravamo pretpostavke režima rada.

Ako pretpostavimo aktivan režim, ulazna kontura daje
\[V_I-R_BI_{RB}-V_{BE1}=0,\qquad I_B=I_{RB}=\frac{V_I-V_{BE1}}{R_B}.\]
Pri $V_{BE1}=V_{BE}=0.6\,\mathrm{V}$ dobijamo negativnu baznu struju. To protivreči zahtevu $I_B\geq0$ u aktivnom režimu. Isti zaključak sledi iz pretpostavke zasićenja, kada bismo koristili $V_{BES}=0.7\,\mathrm{V}$.

Preostaje zakočen režim. Uz $I_B=I_{RB}=0$ ista ulazna jednačina daje $V_{BE1}=V_I=0$, pa je uslov $V_{BE1}<V_{\gamma T}$ zadovoljen. Izlazna kontura, nezavisno od režima, daje
\[V_O=V_{CC}-R_CI_{RC}.\]
Za neopterećeno kolo $I_{RC}=I_C$. Tranzistor je zakočen i $I_{C1}=0$, pa je $V_O=V_{CC}$.
